\chapter{Cadre général du projet et état de l’art}

\section{Introduction}

Ce premier chapitre est consacré à la présentation de l’étude préalable du projet réalisé. Cette présentation commence d’abord par situer le projet dans son cadre général, puis nous présentons l’organisme d’accueil ainsi qu’une étude critique des solutions existantes dans le domaine de l’analyse vidéo sportive. Ensuite, nous décrivons la problématique traitée, la solution proposée et les objectifs visés. Enfin, nous concluons ce chapitre en exposant la méthodologie de travail adoptée ainsi que les technologies utilisées pour la conception et le développement de la plateforme SmartPlay.

\section{Cadre général}

Le projet de fin d’études consiste à réaliser une plateforme intelligente d’analyse vidéo dédiée au padel, baptisée SmartPlay. Cette solution permet d’exploiter des techniques avancées d’intelligence artificielle et de vision par ordinateur afin d’analyser automatiquement les matchs à partir de simples vidéos.

La plateforme est capable de détecter les joueurs, de suivre la balle, d’identifier les frappes, de générer des statistiques détaillées et de proposer des recommandations tactiques aux joueurs et aux entraîneurs. L’objectif est de rendre l’analyse de performance plus accessible, plus rapide et plus précise, aussi bien pour les amateurs que pour les professionnels.
\section{Présentation de l’organisme d’accueil}

\subsection{Présentation de l’entreprise}

Ultima est une startup tunisienne, implantée à Tunis, qui s’est positionnée à l’avant-garde du développement de systèmes intelligents. 
Spécialisée dans l’intégration de solutions d’intelligence artificielle dans le domaine de sport, cette sturtup conçoit et déploie des produits à haute valeur ajoutée pour les secteurs de la santé , de la gestion de 
L’information des sports collectifs et de l’automatisation des processus sportifs.Sa culture d’innovation et son expertise technique en font un environnement dynamique et stimulant.



\begin{figure}[H]
    \centering
    \maybeincludegraphics[width=0.45\textwidth]{images_chap1/logo.png}
    \caption{Logo de l’entreprise Ultima}
    \label{fig:logo_ultima}
\end{figure}
\subsection{
Encadrement du projet}

Le projet SmartPlay a été réalisé sous la supervision d’un encadrant académique et d’un encadrant professionnel, qui ont assuré le suivi technique et méthodologique du travail.

\section{Présentation du projet}

\subsection{Contexte du projet}

Le padel connaît une croissance rapide à l’échelle mondiale. Cette discipline, qui combine des éléments du tennis et du squash, attire un nombre croissant de joueurs amateurs et professionnels.

Cependant, contrairement à d’autres sports comme le football ou le basketball, les outils d’analyse automatisée pour le padel restent peu accessibles. Les analyses sont souvent réalisées manuellement, ce qui nécessite beaucoup de temps et d’expertise.

Dans ce contexte, SmartPlay vise à démocratiser l’analyse de performance en proposant une solution intelligente capable d’extraire automatiquement des informations stratégiques à partir d’une simple vidéo.
\subsection{Problématique}

Le défi central de ce projet réside dans la conception d’un système capable d’analyser automatiquement des matchs de padel à partir de flux vidéo bruts, sans intervention humaine. Cette automatisation soulève plusieurs problématiques techniques interdépendantes.

\paragraph{Détection et suivi des éléments de jeu}
La faible taille de la balle, sa vitesse élevée ainsi que les phénomènes d’occlusion, combinés aux variations d’éclairage et d’angle de capture, rendent particulièrement complexe la détection fiable et le suivi continu des différentes entités du jeu, notamment les joueurs, la balle et les zones du terrain.

\paragraph{Extraction de statistiques et d’indicateurs de performance}
Les données visuelles brutes doivent être transformées en informations analytiques exploitables permettant d’évaluer les déplacements des joueurs, la dynamique des échanges ainsi que les comportements tactiques observés durant le match.

\paragraph{Interprétation cohérente des événements de jeu}
Le système doit être capable d’analyser les interactions spatio-temporelles entre les différentes entités détectées afin d’identifier et d’interpréter correctement les événements significatifs du match. Cette tâche demeure complexe en raison de la nature dynamique du jeu et des incertitudes pouvant apparaître lors des phases de détection et de suivi.

\vspace{0.5cm}

Ces différentes problématiques convergent vers la question centrale suivante :

\begin{quote}
\textbf{Comment concevoir un système intelligent capable d’analyser automatiquement un flux vidéo de match de padel, d’extraire des statistiques de performance avancées ainsi que des indicateurs tactiques fiables, tout en assurant une interprétation cohérente des événements de jeu observés durant le match ?}
\end{quote}
\section{Solution proposée}

Ce projet propose la conception d’un système intelligent d’analyse automatisée des matchs de padel par exploitation de flux vidéo. La solution repose sur une architecture modulaire en pipeline structurée autour de cinq composants complémentaires.

\paragraph{Détection visuelle}
Des modèles d’apprentissage profond assurent la détection des joueurs, de la balle et des zones du terrain dans des scènes dynamiques et visuellement complexes.

\paragraph{Suivi temporel}
Un mécanisme de suivi multi-objets reconstruit les trajectoires des entités détectées de manière continue entre les images successives.

\paragraph{Modélisation géométrique}
Une transformation homographique établit la correspondance entre le référentiel image et le référentiel réel du terrain, permettant une localisation précise des entités dans les zones de jeu.

\paragraph{Analyse événementielle}
Une machine à états finis modélise les phases des échanges et leurs transitions, assurant une interprétation cohérente du déroulement du match.

\paragraph{Génération de statistiques}
Les données interprétées alimentent un module de production de métriques avancées : déplacements, trajectoires, zones d’activité et indicateurs tactiques, exploitables par les joueurs, entraîneurs et analystes.

\vspace{0.5cm}

En définitive, la solution proposée combine vision par ordinateur, analyse spatio-temporelle et modélisation logique afin de transformer des flux vidéo bruts en données analytiques exploitables pour le suivi des performances et l’analyse tactique des matchs de padel.

\subsection{Objectifs du projet}

L’objectif principal de ce projet est de concevoir et de développer une plateforme intelligente capable d’analyser automatiquement les matchs de padel à partir de flux vidéo. Cette solution vise à exploiter les techniques d’intelligence artificielle et de vision par ordinateur afin de fournir des statistiques avancées, des indicateurs de performance pertinents ainsi que des informations exploitables sur le déroulement du jeu. L’objectif est de transformer les données vidéo brutes en connaissances structurées utiles pour les joueurs, les entraîneurs et les analystes sportifs.

Afin d’atteindre cet objectif global, le projet est structuré autour de trois objectifs spécifiques complémentaires :

\textbf{Objectif 1 : Automatisation de l’analyse des matchs de padel}  

Le premier objectif consiste à développer un système capable d’analyser automatiquement les matchs de padel à partir de vidéos. Ce système repose sur des techniques de vision par ordinateur permettant la détection et le suivi des joueurs, ainsi que l’identification de la balle et des événements clés du jeu. L’ensemble de ces informations est ensuite utilisé pour générer des statistiques fiables et des indicateurs de performance précis, reflétant le déroulement réel des matchs.

\vspace{0.3cm}

\textbf{Objectif 2 : Génération de recommandations personnalisées}  

Le deuxième objectif vise à exploiter les données analytiques produites par le système afin de proposer des recommandations personnalisées. Ces recommandations ont pour but d’aider les joueurs et les entraîneurs à mieux comprendre leurs performances, à identifier leurs points faibles et à améliorer leurs stratégies tactiques. Cette approche permet d’apporter une valeur ajoutée en transformant les statistiques en insights exploitables.

\vspace{0.3cm}

\textbf{Objectif 3 : Conception d’une interface intuitive}  

Enfin, le troisième objectif consiste à concevoir et développer une interface utilisateur ergonomique et intuitive. Cette interface permet de visualiser facilement les statistiques générées, de suivre l’analyse des matchs en temps réel et d’interpréter efficacement les données produites par le système. L’objectif est de rendre la plateforme accessible même aux utilisateurs non techniques tout en assurant une expérience fluide et interactive.

\subsection{Défis techniques}

Le développement de SmartPlay présente plusieurs défis techniques liés à l’analyse automatisée des matchs de padel.

\begin{itemize}
    \item \textbf{Petite taille et grande vitesse de la balle ainsi que les réflexions sur les vitres}  
    La petite taille de la balle, sa vitesse élevée ainsi que les réflexions sur les parois vitrées compliquent sa détection et son suivi précis durant les échanges.

    \item \textbf{Occultation fréquente des joueurs et changements d’identification}  
    Les interactions rapprochées et les déplacements rapides des joueurs provoquent des occlusions pouvant entraîner des pertes de suivi ou des erreurs d’identification.

    \item \textbf{Variabilité des conditions d’éclairage et calibration des caméras}  
    Les variations d’éclairage, les différences d’angle de vue et la qualité des vidéos peuvent affecter les performances des modèles de détection et de suivi.

    \item \textbf{Détection précise des frappes et des événements de jeu}  
    La rapidité des échanges rend difficile l’identification fiable des frappes, des rebonds et des autres événements importants du match.
\end{itemize}

\subsection{Solutions existantes}

\subsubsection{Padelytics}
Padelytics est une plateforme d'analyse de matchs de padel proposant la génération automatique de statistiques, de cartes de chaleur (\textit{heatmaps}) et de séquences vidéo synthétiques (\textit{highlights}). Elle s'appuie sur des algorithmes de vision par ordinateur pour détecter les joueurs et suivre la trajectoire de la balle. Son déploiement reste limité aux clubs équipés d'une infrastructure caméra compatible et son modèle tarifaire exclut de fait les structures amateur.

\subsubsection{Hudl}
Hudl est une plateforme multi-sport d'analyse vidéo largement adoptée dans les niveaux semi-professionnels et professionnels. Elle offre des outils de découpage vidéo, d'annotation tactique et de partage entre entraîneurs. Toutefois, son architecture généraliste ne prend pas en charge les spécificités du padel, notamment la détection des frappes et le suivi précis de la balle.

\subsubsection{Dartfish}
Dartfish est un logiciel d'analyse biomécanique et tactique utilisé dans les fédérations sportives nationales et les programmes olympiques. Il permet une analyse image par image et des superpositions comparatives. Son positionnement haut de gamme, sa complexité de prise en main et l'absence de modules dédiés au padel en limitent l'applicabilité dans ce contexte.

\subsection{Analyse comparative des solutions existantes}

Le tableau~\ref{tab:comparaison} présente une évaluation multicritère des solutions identifiées au regard des exigences fonctionnelles et non-fonctionnelles retenues pour SmartPlay.

\begin{table}[H]
\centering
\caption{Comparaison multicritère des solutions existantes et de SmartPlay}
\label{tab:comparaison}
\renewcommand{\arraystretch}{1.3}
\begin{tabular}{|p{4.5cm}|c|c|c|c|}
\hline
\textbf{Critère} & \textbf{Padelytics} & \textbf{Hudl} & \textbf{Dartfish} & \textbf{SmartPlay} \\
\hline
Spécialisation padel & \checkmark & \texttimes & \texttimes & \checkmark \\
\hline
Détection automatique des joueurs & \checkmark & \checkmark & \checkmark & \checkmark \\
\hline
Suivi de la balle & \checkmark & Partiel & Partiel & \checkmark \\
\hline
Détection des frappes & \checkmark & \texttimes & \texttimes & \checkmark \\
\hline
Génération de recommandations tactiques & Limité & \texttimes & \texttimes & \checkmark \\
\hline
Accessibilité aux structures amateur & \texttimes & \texttimes & \texttimes & \checkmark \\
\hline
Architecture modulaire et évolutive & Non documenté & \texttimes & \texttimes & \checkmark \\
\hline
Déploiement sans infrastructure dédiée & \texttimes & Partiel & \texttimes & \checkmark \\
\hline
\end{tabular}
\end{table}

\subsection{Bilan et positionnement de SmartPlay}

L'analyse comparative met en évidence trois lacunes structurelles communes aux solutions examinées.

\textbf{Barrière à l'entrée économique.} Les tarifs pratiqués par Padelytics, Hudl et Dartfish sont calibrés pour des organisations professionnelles ou semi-professionnelles, excluant les clubs amateurs et les académies de formation disposant de budgets opérationnels limités.

\textbf{Déficit de spécialisation disciplinaire.} Hudl et Dartfish reposent sur des modèles generiques non adaptés aux contraintes cinématiques du padel : vitesse et trajectoire de balle spécifiques, configuration fermée du terrain, patterns de jeu en double. L'absence de détection des frappes dans ces deux solutions illustre ce manque d'adaptation.

\textbf{Absence de recommandations tactiques automatisées.} Aucune solution existante ne propose de génération automatique de recommandations tactiques fondées sur l'analyse statistique des séquences de jeu. Padelytics offre un amorce fonctionnelle sur ce point, mais sans module d'inférence documenté ni personnalisation par profil de joueur.

SmartPlay se positionne comme une réponse directe à ces trois lacunes. La solution combine une spécialisation disciplinaire complète pour le padel, une architecture modulaire permettant un déploiement progressif sans infrastructure lourde, et un moteur de recommandations tactiques fondé sur l'apprentissage automatique. Son modèle de déploiement vise explicitement l'accessibilité aux clubs amateurs, segment non couvert par les acteurs actuels du marché.

\section{Méthodologie de travail}

Pour organiser efficacement le développement de la plateforme SmartPlay, nous avons adopté une approche agile basée principalement sur la méthodologie \textbf{Scrum}. Cette méthode nous a permis de structurer le projet en cycles de développement courts appelés \textit{sprints}, chacun ayant des objectifs précis et des livrables clairement définis.

Chaque sprint correspond à une période de travail d’une semaine, au cours de laquelle un ensemble de fonctionnalités est conçu, développé, testé et validé. Cette organisation a permis de progresser de manière incrémentale tout en maintenant une visibilité claire sur l’état d’avancement du projet.

\subsection{Organisation de la méthodologie adoptée}

La méthodologie mise en place repose sur les éléments suivants :

\begin{itemize}
    \item \textbf{Product Backlog} : liste complète des fonctionnalités à développer, regroupant l’ensemble des besoins fonctionnels et techniques du projet.
    
    \item \textbf{Sprint Planning} : réunion hebdomadaire permettant de définir les objectifs du sprint et les tâches prioritaires à réaliser.
    
    \item \textbf{Daily Update} : message quotidien résumant les tâches réalisées, les travaux prévus pour la journée et les éventuels obstacles rencontrés.
    
    \item \textbf{Weekly Meeting} : réunion hebdomadaire avec l’encadrant pour présenter les résultats obtenus, recueillir les remarques et ajuster les priorités.
    
    \item \textbf{Sprint Review et Retrospective} : évaluation des fonctionnalités développées et identification des améliorations à apporter.
\end{itemize}

Cette organisation nous a permis d’assurer un suivi rigoureux de l’avancement du projet, de documenter les progrès réalisés et de valider progressivement les différentes composantes de la solution.

\begin{figure}[H]
    \centering
    \maybeincludegraphics[width=0.85\textwidth]{images_chap1/scrum_cycle.png}
    \caption{Cycle Scrum adopté dans le projet SmartPlay}
    \label{fig:scrum_cycle}
\end{figure}

La figure \ref{fig:scrum_cycle} illustre les principales étapes du cycle Scrum : la définition du Product Backlog, la planification du sprint, le suivi quotidien à travers les Daily Updates, le développement, la revue de sprint et la rétrospective.

\subsection{Daily Update}

Chaque jour, un message de suivi était rédigé afin de communiquer l’état d’avancement du projet. Ce message comportait généralement :

\begin{itemize}
    \item les tâches réalisées la veille ;
    \item les tâches prévues pour la journée ;
    \item les difficultés ou blocages rencontrés.
\end{itemize}

Cette pratique a permis de maintenir une progression continue et d’identifier rapidement les problèmes techniques.

\subsection{Weekly Meeting}

Une réunion hebdomadaire était organisée avec l’encadrant afin de :

\begin{itemize}
    \item présenter les fonctionnalités développées ;
    \item démontrer les résultats obtenus ;
    \item recueillir les remarques et recommandations ;
    \item ajuster les priorités du sprint suivant.
\end{itemize}

Ces réunions ont joué un rôle essentiel dans la validation progressive du projet.

\subsection{Planification des sprints}

Le développement du projet SmartPlay a été structuré en onze sprints successifs, correspondant aux principales étapes de conception et de déploiement de la solution.

\begin{table}[H]
\centering
\caption{Planification des sprints du projet SmartPlay}
\label{tab:sprints}
\renewcommand{\arraystretch}{1.3}
\begin{tabular}{|c|p{5cm}|p{7cm}|}
\hline
\textbf{Sprint} & \textbf{Intitulé} & \textbf{Objectif principal} \\
\hline
1 & Mise en place du projet & Configuration de l’environnement de développement et organisation des données. \\
\hline
2 & Prétraitement vidéo et détection du terrain & Synchronisation des vidéos et calcul de l’homographie. \\
\hline
3 & Détection et suivi & Détection des joueurs avec YOLO, tracking avec ByteTrack et détection de la balle avec TrackNet. \\
\hline
4 & Reconstruction 3D et analyse des événements & Détection des frappes, rebonds et segmentation des échanges. \\
\hline
5 & Classification des coups & Identification des types de coups et analyse de la trajectoire de la balle. \\
\hline
6 & Analyse des performances & Calcul des distances, vitesses et indicateurs de performance. \\
\hline
7 & Génération des statistiques & Production des statistiques détaillées du match. \\
\hline
8 & Modèles prédictifs & Reconnaissance des schémas tactiques et prédiction des coups. \\
\hline
9 & Production vidéo automatisée & Génération automatique des highlights et overlays graphiques. \\
\hline
10 & Plateforme cloud et coaching & Recommandations personnalisées et fonctionnalités sociales. \\
\hline
11 & Déploiement en production & Optimisation, MLOps et intégration avec SUMMA. \\
\hline
\end{tabular}
\end{table}

\begin{figure}[H]
    \centering
    \maybeincludegraphics[width=0.80\textwidth]{images_chap1/sprint_roadmap.png}
    \caption{Roadmap des sprints du projet SmartPlay}
    \label{fig:sprint_roadmap}
\end{figure}

La figure \ref{fig:sprint_roadmap} présente la succession chronologique des différents sprints, depuis la mise en place de l’environnement de développement jusqu’au déploiement final de la solution.

\subsection{Déroulement d’un sprint}

Chaque sprint a été réalisé selon les étapes suivantes :

\begin{enumerate}
    \item Planification des objectifs et des tâches.
    \item Développement des fonctionnalités.
    \item Suivi quotidien à travers les Daily Updates.
    \item Tests et validation.
    \item Revue de sprint.
    \item Rétrospective.
\end{enumerate}

\subsection{Avantages de la méthodologie adoptée}

L’approche Scrum a apporté plusieurs bénéfices :

\begin{itemize}
    \item une planification claire des objectifs ;
    \item un suivi quotidien rigoureux ;
    \item une communication régulière avec l’encadrant ;
    \item une amélioration continue de la solution ;
    \item une intégration progressive des différents modules.
\end{itemize}

En adoptant cette méthodologie agile, nous avons pu structurer efficacement le développement de SmartPlay, suivre l’avancement du projet au quotidien et garantir une progression continue jusqu’à la réalisation d’une solution complète et industrialisable.
\section{Technologies utilisées}

L’architecture technique de SmartPlayAI repose sur une pile technologique moderne combinant la vision par ordinateur, l’intelligence artificielle, l’analyse de données et les technologies de déploiement. Cette stack a été sélectionnée afin de garantir des performances élevées, une bonne modularité ainsi qu’une évolutivité permettant l’intégration future de nouvelles fonctionnalités.

Le développement de la solution a été réalisé principalement avec Python 3.10, un langage particulièrement adapté au prototypage rapide et au développement d’applications d’intelligence artificielle.

\subsection{Interface utilisateur}

L’interface de visualisation a été développée avec Streamlit, un framework permettant de créer rapidement des tableaux de bord interactifs. Cette interface offre aux utilisateurs la possibilité de charger des vidéos, d’exécuter l’analyse automatique et de consulter les statistiques générées sous forme de graphiques et de visualisations dynamiques.

Pour la représentation graphique des données, Plotly a été utilisé afin de produire des courbes interactives, des heatmaps et des tableaux analytiques.

\subsection{Vision par ordinateur}

Le traitement vidéo repose sur OpenCV, qui assure la lecture, le traitement et l’annotation des images. L’outil FFmpeg est utilisé pour la conversion, la compression et la manipulation efficace des vidéos haute résolution.

La bibliothèque Supervision (développée par Roboflow) facilite la gestion des détections, la définition de zones polygonales correspondant au terrain de padel et l’annotation visuelle des résultats.

\subsection{Analyse de données}

Les calculs numériques, notamment les positions, vitesses et accélérations des joueurs et de la balle, sont effectués grâce à NumPy et Pandas. Ces bibliothèques permettent de structurer les données et de produire des indicateurs statistiques détaillés.

\subsection{Intelligence artificielle}

La solution intègre plusieurs modèles et algorithmes d’intelligence artificielle :

\begin{itemize}
    \item YOLO pour la détection des joueurs et des éléments du terrain ;
    \item TrackNet pour le suivi précis de la balle ;
    \item Des trackers personnalisés pour le suivi des joueurs, de la balle et des points clés du terrain ;
    \item Des algorithmes d’estimation de vitesse et de projection spatiale sur une représentation 2D du court.
\end{itemize}

\subsection{Infrastructure et déploiement}

Afin de garantir la portabilité de la solution et l’isolation de ses dépendances, l’application est conteneurisée avec Docker. L’orchestration des différents services est assurée par Docker Compose, ce qui simplifie considérablement le déploiement sur différentes machines.

\subsection{Fiche technique synthétique}

\begin{table}[H]
\centering
\caption{Fiche technique de SmartPlay}
\begin{tabular}{|l|l|}
\hline
\textbf{Composant} & \textbf{Technologies} \\
\hline
Langage principal & Python 3.10 \\
Interface utilisateur & Streamlit, Plotly \\
Vision par ordinateur & OpenCV, Supervision, FFmpeg \\
Analyse de données & Pandas, NumPy \\
Intelligence artificielle & YOLO, TrackNet, trackers personnalisés \\
Infrastructure & Docker, Docker Compose \\
\hline
\end{tabular}
\end{table}

\section{Conclusion}

Dans ce chapitre, nous avons présenté le cadre général du projet SmartPlay, son contexte, la problématique traitée ainsi que les objectifs visés. Nous avons également étudié les solutions existantes, mis en évidence leurs limites et positionné notre solution par rapport à ces outils.

Enfin, nous avons décrit la méthodologie de travail adoptée ainsi que les principales technologies utilisées pour la conception de la plateforme.

Le chapitre suivant sera consacré à l’analyse des besoins et à la conception détaillée de la solution.
